\chapter{Washing hands}
\index{Washing hands}

\section*{Definition and Overview}
Hand hygiene, universally recognized as hand washing or hand sanitisation, is the practice of cleaning one's hands with soap and water or using an antiseptic agent (such as an alcohol-based hand rub) to remove dirt, organic material, and microorganisms. In clinical medicine and public health, hand hygiene is documented as the single most critical, cost-effective, and practical intervention to reduce the transmission of infectious pathogens and prevent healthcare-associated infections (HAIs). The hands are the primary vectors of contact transmission; they interface continuously with the external environment, touch high-frequency contact surfaces, and interact with mucosal membranes, open wounds, and medical devices. Consequently, hand hygiene serves as a fundamental barrier that interrupts the chain of infection.

Historically, the clinical significance of hand hygiene was established in the mid-nineteenth century. In 1847, the Hungarian physician Ignaz Semmelweis observed that maternal mortality due to puerperal fever (childbed fever) was drastically higher in obstetric clinics where births were attended by medical students who had recently performed autopsies compared to clinics run by midwives. Semmelweis proposed that students carried ``cadaveric particles'' on their hands and mandated that they wash their hands in a chlorinated lime solution before examining patients. This simple intervention caused maternal mortality rates to plunge from approximately 18\% to under 2\%. Shortly thereafter, Florence Nightingale implemented rigorous hand washing and sanitation protocols during the Crimean War, resulting in a dramatic decline in soldier mortality from infectious diseases. In the modern era, hand hygiene forms the foundation of standard precautions in healthcare settings, guided by global initiatives such as the World Health Organization's (WHO) ``Save Lives: Clean Your Hands'' campaign and the Centers for Disease Control and Prevention's (CDC) clean hands guidelines.

\section*{Epidemiology}
The epidemiology of hand hygiene compliance and its impact on disease transmission reveals a stark contrast between its simple implementation and its critical importance in global health. In healthcare facilities worldwide, healthcare-associated infections represent a major patient safety concern. At any given time, out of every 100 hospitalized patients, approximately seven in high-income countries and fifteen in low- and middle-income countries will acquire at least one HAI. Despite the clear correlation between clean hands and infection control, observational studies consistently demonstrate that baseline compliance with hand hygiene protocols among healthcare workers is remarkably low, frequently averaging less than 40\% without active educational and institutional interventions. Compliance rates vary significantly by profession, with nursing staff generally demonstrating higher adherence than medical physicians, and by clinical setting, where compliance is paradoxically lower in high-stress, high-workload environments such as intensive care units (ICUs) compared to general medical wards.

In the community, inadequate hand hygiene is a primary driver of endemic and epidemic outbreaks of gastrointestinal and respiratory illnesses. Epidemiological data indicates that hand washing with soap and water can reduce diarrheal disease morbidity by approximately 30\% to 48\%, which is particularly significant in developing nations where diarrheal diseases remain a leading cause of death in children under the age of five. Furthermore, systematic reviews indicate that proper hand hygiene reduces the risk of acute respiratory tract infections, including seasonal influenza, rhinovirus, and coronavirus infections, by 20\% to 23\%. Demographic analyses of community hand washing habits show that adherence is statistically lower in males than in females after using public restrooms. Compliance also correlates closely with socioeconomic factors; populations lacking reliable access to clean running water, sanitation infrastructure, and basic hygiene education experience disproportionately higher rates of hand-transmitted infectious diseases.

\section*{Aetiology and Pathophysiology}
The biological basis of hand hygiene involves understanding the microbiological composition of the skin, the mechanisms of pathogen transfer, and the physiological effects of hand hygiene agents on the skin barrier. The aetiological and pathophysiological principles include:
\begin{itemize}
    \item \textbf{Transient Skin Flora}: These microorganisms colonize the superficial layers of the skin (the stratum corneum) and are acquired during direct contact with patients, contaminated surfaces, or environmental reservoirs. Unlike resident flora, transient microbes do not multiply on the skin and are easily removed by mechanical washing. Key transient pathogens include Gram-negative bacilli such as \textit{Escherichia coli}, \textit{Klebsiella pneumoniae}, and \textit{Pseudomonas aeruginosa}, as well as Gram-positive bacteria like \textit{Staphylococcus aureus} and various species of \textit{Streptococcus}. Non-enveloped viruses such as Norovirus and Rotovirus are also critical transient pathogens.
    \item \textbf{Resident Skin Flora}: These are permanent microbiological inhabitants of the deeper layers of the skin, including the stratum corneum crevices, hair follicles, and sebaceous glands. They are generally of low virulence and include coagulase-negative staphylococci (such as \textit{Staphylococcus epidermidis}), \textit{Micrococcus} species, and \textit{Corynebacterium} species. Resident flora are highly resistant to mechanical washing but can be temporarily reduced by chemical antiseptics. They play a protective role by preventing colonization by more pathogenic transient species, though they can cause infection if introduced into sterile body sites via invasive procedures.
    \item \textbf{Contact Transmission Pathway}: The hands facilitate indirect or direct contact transmission. In direct transmission, pathogens are transferred directly from an infected individual's skin or mucous membranes to the hands of another person. In indirect transmission, a contaminated hand touches an inanimate object, known as a fomite (such as bed rails, stethoscopes, doorknobs, or cell phones), depositing pathogens that can survive for hours or days. The next person to touch the fomite picks up these pathogens on their hands, subsequently introducing them into their own body by touching their eyes, nose, or mouth, or by handling food.
    \item \textbf{Mechanism of Surfactant Action}: Soap is a surfactant containing molecules with a hydrophilic (water-attracting) head and a hydrophobic (water-repelling) tail. When soap is applied to hands and scrubbed, the hydrophobic tails bind to the lipids in the cell membranes of bacteria and the envelopes of viruses, as well as to dirt and grease. The hydrophilic heads align toward the water. This forms spherical structures called micelles that trap the lipids, dirt, and pathogens. The mechanical friction of rubbing the hands together dislodges these micelles from the skin surface, allowing them to be rinsed away in the water stream.
    \item \textbf{Epidermal Barrier Disruption}: While hand hygiene is vital for infection control, the pathophysiology of frequent hand washing can lead to skin damage. The stratum corneum acts as a physical barrier composed of corneocytes surrounded by an intercellular lipid matrix containing ceramides, cholesterol, and free fatty acids. Repeated exposure to anionic surfactants in soaps and hot water emulsifies and strips away these essential lipids. This causes denaturation of keratin proteins, leading to a rise in transepidermal water loss (TEWL). The subsequent cellular dehydration triggers a pro-inflammatory cytokine cascade, leading to the development of irritant contact dermatitis.
\end{itemize}

\section*{Clinical Features}
The clinical features associated with hand hygiene must be evaluated from two distinct perspectives: the presentation of infections that result from inadequate hand hygiene compliance, and the dermatological manifestations resulting from excessive or improper hand washing practices.

When hand hygiene is deficient, the resulting clinical presentations vary according to the transmitted pathogen:
\begin{itemize}
    \item \textbf{Gastrointestinal Infections}: Often presenting as acute gastroenteritis, patients exhibit symptoms such as sudden onset of nausea, projectile vomiting, watery or bloody diarrhea, diffuse abdominal cramping, and low-grade to high-grade fever. Dehydration is a common clinical sign, manifested by dry mucous membranes, delayed capillary refill, tachycardia, and orthostatic hypotension.
    \item \textbf{Respiratory Infections}: Transmission of respiratory pathogens yields clinical signs such as rhinorrhea, nasal congestion, pharyngitis, productive or non-productive cough, dyspnea, wheezing, and systemic symptoms including myalgias, arthralgias, headache, and severe fatigue.
    \item \textbf{Surgical Site and Wound Infections}: Characterized by localized erythema, warmth, edema, purulent exudate, and progressive tenderness at the site of surgical incisions or open wounds, often accompanied by systemic leukocytosis and fever.
\end{itemize}

Conversely, the clinical features of hand-hygiene-induced skin damage (irritant contact dermatitis) are highly prevalent among healthcare workers and individuals performing frequent wet work:
\begin{itemize}
    \item \textbf{Mild to Moderate Xerosis}: The initial presentation is characterized by dryness, roughness, and a dull appearance of the skin, predominantly affecting the dorsum of the hands and the interdigital webs.
    \item \textbf{Erythema and Scaling}: As inflammation progresses, the skin becomes erythematous, with visible scaling, flaking, and superficial peeling.
    Pruritus and Burning: Patients report persistent itching, burning, or stinging sensations, which are typically exacerbated by subsequent hand washing or the application of alcohol-based hand rubs.
    \item \textbf{Fissuring and Lichenification}: Chronic exposure leads to deep, painful linear cracks (fissures) in the skin, particularly over the joints and fingertips, which may bleed. Over time, the skin undergoes lichenification (thickening and hardening of the epidermis with exaggerated skin markings) due to chronic scratching and rubbing.
\end{itemize}

\section*{Differential Diagnosis}
When evaluating a patient presenting with hand skin pathology secondary to frequent hand washing, clinicians must distinguish irritant contact dermatitis from other dermatological conditions. The differential diagnoses include:
\begin{itemize}
    \item \textbf{Allergic Contact Dermatitis}: A type IV delayed-type hypersensitivity reaction to specific allergens present in hand hygiene products, such as fragrances, preservatives (e.g., methylisothiazolinone, parabens), or surfactants (e.g., cocamidopropyl betaine). Unlike irritant contact dermatitis, which is limited to the areas of direct contact, allergic contact dermatitis can spread beyond the exposure site, presents with intense pruritus, and may feature acute vesiculation, oozing, and edema.
    \item \textbf{Atopic Dermatitis}: A chronic, relapsing inflammatory skin disease associated with a personal or family history of atopic diseases (asthma, allergic rhinitis, atopic eczema). It presents as intensely pruritic, erythematous plaques that frequently involve the flexural surfaces of the body, but can present as hand eczema that is highly susceptible to irritation from routine hand washing.
    \item \textbf{Dyshidrotic Eczema}: Also known as pompholyx, this condition is characterized by the sudden eruption of highly pruritic, deep-seated, vesicles on the lateral and palmar aspects of the fingers and palms. These vesicles have a classic ``tapioca pudding'' appearance and resolve with desquamation. While triggered by stress or sweating, it is exacerbated by frequent hand washing.
    \item \textbf{Tinea Manuum}: A superficial dermatophyte infection of the hands, most commonly caused by \textit{Trichophyton rubrum}. It typically presents with a slowly progressive, unilateral, scaling, erythematous plaque with a well-demarcated active border. It is frequently associated with concurrent tinea pedis, presenting as the classic ``two feet, one hand'' syndrome.
    \item \textbf{Palmoplantar Psoriasis}: An autoimmune inflammatory disorder presenting with well-demarcated, thick, erythematous plaques covered by silvery-white scales on the palms. Fissuring is common and painful. Diagnosis is supported by nail changes (pitting, onycholysis) and psoriatic plaques on other extensor surfaces (elbows, knees).
    \item \textbf{Scabies}: An infestation by the mite \textit{Sarcoptes scabiei}, presenting with generalized pruritus that worsens at night. Clinical signs include linear burrows in the web spaces of the fingers, wrist creases, and palms, often accompanied by secondary excoriations and papules.
\end{itemize}

\section*{When to Seek Medical Care}
While standard hand washing is a daily preventive measure, clinical complications arising from poor hand hygiene (infections) or excessive hand washing (severe dermatitis) require prompt medical evaluation.

Individuals should contact a primary care physician or visit an urgent care clinic if they experience:
\begin{itemize}
    \item Severe dryness and erythema of the hands that does not improve within two weeks of using over-the-counter barrier creams and emollients.
    \item Deep, painful fissures on the hands that bleed or interfere with daily activities and professional duties.
    \item Signs of a localized secondary bacterial infection, such as increased warmth, swelling, throbbing pain, pus or purulent drainage, or red streaks spreading up the arm.
    \item Persistent diarrhea or respiratory symptoms lasting more than a few days, especially if accompanied by a high fever.
\end{itemize}

Immediate emergency medical attention must be sought by calling emergency service numbers, such as \textbf{local emergency services} in many countries, \textbf{911} in the United States, or \textbf{999} in the United Kingdom, if the individual displays:
\begin{itemize}
    \item Signs of systemic bacteremia or sepsis, including a high fever (>38.5$^\circ$C or 101.3$^\circ$F) or hypothermia (<36$^\circ$C or 96.8$^\circ$F), rapid breathing, tachycardia, severe chills, shivering, confusion, or extreme lethargy.
    \item Severe dehydration from gastroenteritis, indicated by an inability to keep fluids down, absence of urination for over 8 hours, extreme dizziness, or loss of consciousness.
    \item Severe respiratory distress, such as stridor, grunting, intercostal retractions, cyanosis (blue-tinted lips or fingernails), or an inability to speak in full sentences.
\end{itemize}

\section*{Diagnosis and Investigations}
In clinical practice, the diagnosis of hand-hygiene-related issues encompasses auditing hand hygiene compliance, conducting epidemiological investigations during outbreaks, and diagnosing dermatological side effects.
\begin{itemize}
    \item \textbf{Direct Observation Audits}: This is the primary method to assess hand hygiene compliance in healthcare settings. Trained observers systematically monitor healthcare workers' compliance with the WHO ``My 5 Moments for Hand Hygiene'' framework. These moments include: (1) before touching a patient, (2) before a clean/aseptic procedure, (3) after body fluid exposure risk, (4) after touching a patient, and (5) after touching patient surroundings.
    \item \textbf{Indirect Monitoring}: Compliance can also be audited by measuring product consumption, which involves tracking the volume of soap and alcohol-based hand rubs used per 1,000 patient-days. Additionally, electronic monitoring systems utilize sensors on sanitiser dispensers and healthcare workers' badges to record hand hygiene opportunities and actions in real-time.
    \item \textbf{Microbiological Surveillance}: During suspected hospital outbreaks of multidrug-resistant organisms, microbiological hand cultures may be performed. Techniques include the ``fingerprint impression'' method, where fingers are pressed directly onto selective agar plates, or the ``glove juice'' technique, where the hand is massaged inside a sterile glove containing buffer solution, which is then cultured to quantify transient pathogens.
    \item \textbf{Adenosine Triphosphate (ATP) Bioluminescence Assays}: Used to evaluate the cleanliness of hands or surfaces. A swab of the skin is reacted with a luciferase enzyme reagent. The level of light produced is proportional to the amount of biological material (ATP) present, providing rapid, objective feedback on hand hygiene effectiveness.
    \item \textbf{Dermatological Evaluation}: When managing hand dermatitis, a thorough history is obtained regarding occupational exposures, wet work duration, soap types, and personal history of atopy. Physical examination evaluates the distribution and morphology of the lesions.
    \item \textbf{Patch Testing}: If allergic contact dermatitis is suspected, epicutaneous patch testing is performed. Standardized allergens are applied to the patient's back under occlusion for 48 hours, and read at 48 and 72 to 96 hours to identify specific causative allergens, such as preservatives or fragrances in soaps.
\end{itemize}

\section*{Management}
The clinical management of hand hygiene involves the proper implementation of sanitation techniques, institutional guidelines, and the medical management of hand-washing-related skin pathology.

The standard techniques for hand hygiene are divided into two main modalities:
\begin{itemize}
    \item \textbf{Hand Washing with Soap and Water}: Indicated when hands are visibly dirty, contaminated with proteinaceous material, soiled with blood or other body fluids, or after using the restroom. The entire procedure should take 40 to 60 seconds:
    \begin{enumerate}
        \item Wet hands with warm water and apply a sufficient quantity of soap to cover all hand surfaces.
        \item Rub hands palm to palm, then rub the right palm over the left dorsum with interlaced fingers and vice versa.
        \item Rub palm to palm with fingers interlaced, and rub the backs of fingers to opposing palms with fingers interlocked.
        \item Perform rotational rubbing of the left thumb clasped in the right palm and vice versa.
        \item Rub the fingertips of the right hand rotationally in the left palm and vice versa.
        \item Rinse hands thoroughly with water and dry them completely with a single-use disposable paper towel. Use the towel to turn off the faucet to prevent re-contamination.
    \end{enumerate}
    \item \textbf{Alcohol-Based Hand Rub (ABHR)}: The preferred method for routine hand hygiene in clinical settings when hands are not visibly soiled. ABHR is more effective at reducing bacterial counts, requires less time (20 to 30 seconds), and is less damaging to the skin than soap and water. The procedure involves:
    \begin{enumerate}
        \item Apply a palmful of the alcohol-based formulation into a cupped hand.
        \item Rub hands together, covering all surfaces of the hands and fingers using the same rotational and interlocking steps described for hand washing.
        \item Continue rubbing until the alcohol has completely evaporated and the hands are dry. Do not wipe or rinse the product off.
    \end{enumerate}
\end{itemize}

Medical management of hand dermatitis resulting from frequent hand hygiene includes:
\begin{itemize}
    \item \textbf{Barrier Restoration and Emollients}: Liberal, frequent application of lipid-rich emollients (ointments or heavy creams containing ceramides, dimethicone, petrolatum, or urea) immediately after every hand washing or sanitising cycle. This helps to lock in moisture and rebuild the stratum corneum lipid bilayer.
    \item \textbf{Pharmacotherapy}: For moderate to severe inflammation, topical corticosteroids are prescribed. A medium-potency corticosteroid (such as triamcinolone acetonide 0.1\% or methylprednisolone aceponate) is applied to the affected areas once or twice daily for a limited course of 1 to 2 weeks to reduce erythema, scaling, and pruritus. For secondary bacterial infections, topical antibiotics (e.g., mupirocin 2\% ointment) or oral antibiotics (e.g., cephalexin) are indicated.
    \item \textbf{Product Substitution}: Replacing harsh, high-pH alkaline soaps with mild, fragrance-free synthetic detergent bars (syndets) that have a pH close to the skin's natural pH (5.5). In healthcare settings, utilizing ABHRs that contain built-in emollients (like glycerol) helps mitigate the drying effects of alcohol.
\end{itemize}

\section*{Complications}
The complications associated with hand hygiene occur either as a consequence of non-compliance (pathogen transmission) or as a side effect of excessive, improper hand washing (tissue damage).
\begin{itemize}
    \item \textbf{Healthcare-Associated Infections (HAIs)}: Inadequate hand hygiene among healthcare workers is the primary driver of nosocomial outbreaks. These infections include central line-associated bloodstream infections (CLABSIs), catheter-associated urinary tract infections (CAUTIs), ventilator-associated pneumonia (VAP), and surgical site infections (SSIs), resulting in prolonged hospital stays, increased medical costs, and significant patient mortality.
    \item \textbf{Transmission of Multidrug-Resistant Organisms (MDROs)}: Poor compliance leads to the cross-transmission of highly resistant bacterial strains, such as Methicillin-Resistant \textit{Staphylococcus aureus} (MRSA), Vancomycin-Resistant \textit{Enterococcus} (VRE), and Carbapenem-Resistant Enterobacteriaceae (CRE), limiting therapeutic options for vulnerable patients.
    \item \textbf{Clostridioides difficile Spore Dissemination}: The spores of \textit{C. difficile} are highly resistant to alcohol-based hand rubs. Inadequate mechanical washing with soap and water fails to physically remove these spores from the hands. This allows the pathogen to spread between patients, causing severe, recurrent pseudomembranous colitis.
    \item \textbf{Severe Irritant Contact Dermatitis}: Chronic skin barrier degradation from frequent washing can lead to extensive fissuring and painful lichenification. This compromises the skin's natural defense mechanism, transforming the hands from a tool of patient care into a potential reservoir for pathogenic bacteria.
    \item \textbf{Secondary Bacterial Infections}: Open fissures and excoriations on damaged skin serve as portals of entry for pyogenic bacteria, leading to complications such as acute paronychia, cellulitis, erysipelas, or localized abscesses.
    \item \textbf{Occupational and Psychological Impairment}: Severe hand dermatitis can cause significant pain and physical limitation, preventing healthcare workers and food handlers from performing their duties. This leads to increased sick leave, career transitions, and psychological distress, including anxiety and depression related to chronic discomfort.
\end{itemize}

\section*{Prognosis and Outcomes}
The prognosis of conditions related to hand hygiene is generally excellent, provided appropriate behavioral, clinical, and administrative interventions are implemented.

When institutional compliance with hand hygiene protocols is improved, the clinical outcomes are highly favorable. Data demonstrates that implementing multimodal hand hygiene improvement strategies (e.g., the WHO Multimodal Hand Hygiene Improvement Strategy) leads to a significant and sustained reduction in HAI rates, often decreasing infection incidence by up to 50\%. This reduction is associated with decreased patient morbidity, shorter hospitalizations, lower healthcare expenditures, and decreased transmission of multidrug-resistant pathogens.

For individuals suffering from hand-washing-induced irritant contact dermatitis, the prognosis is highly favorable with proper dermatological management. Most acute cases of irritant dermatitis resolve within 2 to 3 weeks once the offending agents are removed, mild syndet cleansers are substituted for harsh soaps, and lipid-rich emollients are applied consistently. However, in individuals with a genetic predisposition to skin barrier dysfunction (such as atopic dermatitis) or those with ongoing, unavoidable occupational exposure to wet work, the condition can become chronic. In these cases, the prognosis depends on the consistent use of barrier creams, adherence to skin care regimens, and occasional use of topical anti-inflammatory medications. Failure to manage chronic dermatitis can lead to persistent skin barrier impairment, increasing the long-term risk of secondary infections and occupational disability.

\section*{Prevention and Self-Care}
Preventing hand-hygiene-related complications requires balancing effective pathogen removal with the maintenance of skin barrier integrity. The following self-care and preventive strategies are recommended:
\begin{itemize}
    \item \textbf{Optimal Water Temperature}: Always wash hands with lukewarm water rather than hot water. Hot water accelerates the dissolution of the intercellular lipid matrix in the stratum corneum, exacerbating skin dryness and irritation.
    \item \textbf{Gentle Drying Technique}: Avoid rubbing the hands vigorously with paper towels. Instead, gently pat the skin dry to minimize mechanical friction and micro-abrasions, paying close attention to the interdigital web spaces where moisture can collect and cause maceration.
    \item \textbf{Immediate Post-Wash Emolliency}: Apply a fragrance-free, lipid-rich emollient cream or ointment immediately after every hand washing cycle while the skin is still slightly damp. This traps moisture within the stratum corneum and helps restore the epidermal barrier.
    \item \textbf{Appropriate Cleanse Selection}: In the home and community, select mild, fragrance-free liquid cleansers or synthetic detergent bars (syndets) over highly alkaline, heavily perfumed bar soaps. Avoid antibacterial soaps containing harsh active agents like triclosan unless clinically indicated, as they offer no additional benefit over plain soap and increase skin irritation.
    \item \textbf{Preferential Use of Alcohol-Based Hand Rubs}: In clinical settings or when hands are not visibly soiled, prefer alcohol-based hand rubs over soap and water. Modern ABHR formulations contain humectants (e.g., glycerin) that are less damaging to the skin barrier than repeated surfactant exposure.
    \item \textbf{Gloves for Protective Wet Work}: Wear appropriate protective gloves (such as vinyl or nitrile gloves with a cotton liner) when performing household cleaning, dishwashing, or handling chemicals to prevent direct contact with irritants and minimize wet-dry cycles.
    \item \textbf{Nail Hygiene Maintenance}: Keep fingernails short and neatly trimmed. The subungual space (the area beneath the fingernail) is a major reservoir for transient microorganisms and is difficult to clean thoroughly. Avoid artificial nails or chipped nail polish, which can harbor high densities of fungi and bacteria.
    \item \textbf{Self-Monitoring}: Conduct daily inspections of the hands for early signs of skin irritation, such as redness, roughness, or mild itching. Promptly initiate intensive emollient therapy at the earliest sign of dryness to prevent the progression to painful fissures.
\end{itemize}

\section*{Sources and Further Reading}
\begin{itemize}
    \item \textbf{World Health Organization (WHO)}: WHO Guidelines on Hand Hygiene in Health Care. Geneva: World Health Organization; 2009. Available at: \texttt{https://www.who.int/publications/i/item/9789241597906}
    \item \textbf{Centers for Disease Control and Prevention (CDC)}: Guideline for Hand Hygiene in Health-Care Settings. MMWR Recommendations and Reports; 2002. Available at: \texttt{https://www.cdc.gov/mmwr/preview/mmwrhtml/rr5116a1.htm}
    \item \textbf{National Health and Medical Research Council (NHMRC)}: Australian Guidelines for the Prevention and Control of Infection in Healthcare. Canberra: NHMRC; 2019. Available at: \texttt{https://www.nhmrc.gov.au/about-us/publications/australian-guidelines-prevention-and-control-infection-healthcare}
    \item \textbf{National Institutes of Health (NIH)}: MedlinePlus Medical Encyclopedia: Handwashing - Clean Hands Save Lives. Bethesda: National Library of Medicine; 2023. Available at: \texttt{https://medlineplus.gov/handwashing.html}
    \item \textbf{Mayo Clinic}: Handwashing: Do's and Don'ts for Clean Hands. Rochester: Mayo Foundation for Medical Education and Research; 2023. Available at: \texttt{https://www.mayoclinic.org/healthy-lifestyle/adult-health/in-depth/handwashing/art-20046253}
    \item \textbf{Association for Professionals in Infection Control and Epidemiology (APIC)}: Hand Hygiene Resources for Infection Prevention. Arlington: APIC; 2022. Available at: \texttt{https://apic.org/resources/topic-specific-infection-prevention-resources/hand-hygiene/}
\end{itemize}